%%%%%%%%%%%%%%%%%%%%%%%%%%%%%%%%%%%%%%%%%%%%%%%%%%%%%%%%%%%%%%%%%%%%%%%%%%%%%%%%
%% MemCatalog: A DBMS-style Memory Management System for LLM Agents
%% ACM SIGMOD submission (acmart sigconf)
%%%%%%%%%%%%%%%%%%%%%%%%%%%%%%%%%%%%%%%%%%%%%%%%%%%%%%%%%%%%%%%%%%%%%%%%%%%%%%%%

%% For submission and double-blind review, use:
%%   \documentclass[sigconf, anonymous, review]{acmart}
%% For camera-ready, use:
%%   \documentclass[sigconf]{acmart}
\documentclass[sigconf, review, anonymous]{acmart}

\AtBeginDocument{%
  \providecommand\BibTeX{{Bib\TeX}}}

%% Rights management information. Replace with the values supplied by ACM
%% after the rights form is completed.
\setcopyright{acmlicensed}
\copyrightyear{2026}
\acmYear{2026}
\acmDOI{XXXXXXX.XXXXXXX}
\acmConference[SIGMOD '26]{Proceedings of the 2026 International Conference on
  Management of Data}{June 14--19, 2026}{Bengaluru, India}
\acmISBN{978-1-4503-XXXX-X/2026/06}

%% Paper-specific commands, colors, and macros.
%%%%%%%%%%%%%%%%%%%%%%%%%%%%%%%%%%%%%%%%%%%%%%%%%%%%%%%%%%%%%%%%%%%%%%%%%%%%%%%%
% MemCatalog paper commands (ACM SIGMOD / acmart sigconf)
%%%%%%%%%%%%%%%%%%%%%%%%%%%%%%%%%%%%%%%%%%%%%%%%%%%%%%%%%%%%%%%%%%%%%%%%%%%%%%%%

\usepackage{amsthm}
\newtheorem{example}{Example}

\usepackage{xcolor}
\usepackage{multirow}
\usepackage{booktabs}
\usepackage{colortbl}
\usepackage{tikz}
\usepackage{subcaption}
\usepackage{enumitem}
\usepackage{xspace}
\usepackage{array}
\usepackage{makecell}
\usepackage{graphicx}

\usetikzlibrary{shapes.geometric,shapes.symbols,positioning,arrows.meta,fit,backgrounds,calc,shadows.blur}

% ===== Math operators =====
\DeclareMathOperator*{\argmax}{arg\,max}

% ===== System name =====
\newcommand{\sys}{MemCatalog\xspace}
\newcommand{\syst}{MemCatalog\xspace}
\newcommand{\nlsql}{{NL2SQL}\xspace}

% ===== Latin abbreviations =====
\newcommand{\eg}{{\em e.g.,}\xspace}
\newcommand{\ie}{{\em i.e.,}\xspace}
\newcommand{\etc}{{\em etc.}\xspace}
\newcommand{\etal}{{\em et al.}\xspace}

% ===== Paragraph titles (ACM run-in heading convention) =====
\newcommand{\stitle}[1]{\smallskip\noindent\textbf{#1}}
\newcommand{\etitle}[1]{\smallskip\noindent\textit{#1}}

% ===== Color palette =====
\definecolor{cadmiumgreen}{rgb}{0.0, 0.42, 0.24}   % positive / gains
\definecolor{dropred}{rgb}{0.75, 0.22, 0.17}       % negative / drops
\definecolor{softblue}{HTML}{2C5F8D}               % primary highlight
\definecolor{softorange}{HTML}{D97F1F}             % secondary highlight
\definecolor{shadecolor}{RGB}{240,240,240}         % subtle background
\definecolor{darkgray_x11}{HTML}{2E2E2E}
\definecolor{lightblue}{HTML}{E8F0F8}
\definecolor{lightgreen}{HTML}{E8F5E9}
\definecolor{lightorange}{HTML}{FFF3E0}

% ===== Result-cell helpers =====
\newcommand{\best}[1]{\textbf{#1}}
\newcommand{\second}[1]{\underline{#1}}
\newcommand{\gain}[1]{{\color{cadmiumgreen}\small(+#1)}}
\newcommand{\loss}[1]{{\color{dropred}\small(-#1)}}

% ===== Named operator badges =====
\newcommand{\opPOINT}{\textcolor{softblue}{\textsc{point}}}
\newcommand{\opAGG}{\textcolor{softorange}{\textsc{aggregate}}}
\newcommand{\opCMP}{\textcolor{cadmiumgreen}{\textsc{compare}}}

% ===== Benchmark name macros =====
\newcommand{\lmev}{{LME-v2}\xspace}
\newcommand{\lmeo}{{LME-o}\xspace}
\newcommand{\lmes}{{LME-s}\xspace}
\newcommand{\lmem}{{LME-m}\xspace}
\newcommand{\locomo}{{LoCoMo}\xspace}
\newcommand{\evomem}{{EvoMem}\xspace}
\newcommand{\membench}{{MemBench}\xspace}
\newcommand{\dynamicmem}{{DynamicMem}\xspace}

% ===== Query class shorthands =====
\newcommand{\point}{\textsc{point}\xspace}
\newcommand{\agg}{\textsc{aggregate}\xspace}
\newcommand{\cmp}{\textsc{compare}\xspace}

% ===== Global table typography (ACM booktabs style) =====
\renewcommand{\arraystretch}{1.1}
\setlength{\tabcolsep}{4.5pt}

% Group header rows spanning a booktabs table.
\newcommand{\tblgroup}[2]{\multicolumn{#1}{l}{\cellcolor{gray!12}\textit{#2}}}
\newcommand{\ourgroup}[2]{\multicolumn{#1}{l}{\cellcolor{lightblue}\textit{#2}}}

% Allow LaTeX to stretch line ends slightly to avoid marginal overfull boxes.
\setlength{\emergencystretch}{3em}

% ===== Experiment numbering and finding boxes (reference-paper convention) =====
\usepackage{framed}
\newcounter{mycounterexp}
\newcommand{\expnum}{\stepcounter{mycounterexp}\arabic{mycounterexp}}
\newcounter{mycounterfinding}
\newcommand{\findingnum}{\stepcounter{mycounterfinding}\arabic{mycounterfinding}}
\definecolor{shadecolor}{RGB}{242,242,242}
\newcommand{\finding}[1]{%
  \begin{shaded}\noindent\textbf{Finding \findingnum: #1}\end{shaded}}


\begin{document}

\title{\sys{}: A DBMS-style Memory Management System for LLM Agents}

\author{Anonymous Author(s)}
\affiliation{%
  \institution{Anonymous Institution}
  \city{Anonymous City}
  \country{Anonymous Country}
}
\email{anon@anon.org}

\renewcommand{\shortauthors}{Anonymous Author(s)}

\begin{abstract}
LLM agents accumulate memory that grows to hundreds of thousands of tokens per user, and in some settings to millions of tokens. The dominant treatment of this memory remains an unmanaged string collection: raw sessions concatenated into an LLM prompt, or dense vectors retrieved by static similarity. This is the pre-DBMS era of agent memory, in which systems can \emph{store} but cannot \emph{manage} their memory. We introduce \sys{}, a DBMS-style memory management system for LLM agents that lifts agent memory into first-class managed data. \sys{} comprises a foundational framework and two novel optimizations built upon it. (i)~A full query-optimization stack, consisting of a catalog, a zero-supervision planner, and three physical operators (index scan, hash-aggregate, and sort-merge), dispatches each query to the matching physical plan. (ii)~A \emph{Second-Order Memory} logs per-user access patterns and reshapes catalog priorities, a workload-to-storage feedback loop transferred from classical DBMS self-tuning to agent memory for the first time. (iii)~A \emph{Cost-Aware Executor} allocates the LLM call budget by plan class and repurposes self-consistency as a cost signal, not a quality voter. Across eight long-term memory benchmarks, \sys{} matches or improves the strongest prior memory system on every one, by an average of $2.6$ points and by up to $5.6$ points on \lmes. On \membench{} it further reduces LLM calls by $50.7\%$ relative to a Uniform-$K{=}3$ reference while attaining higher accuracy.
\end{abstract}

%% CCS concepts. Generated from https://dl.acm.org/ccs.
\begin{CCSXML}
<ccs2012>
   <concept>
       <concept_id>10002951.10002952.10002953</concept_id>
       <concept_desc>Information systems~Data management systems</concept_desc>
       <concept_significance>500</concept_significance>
       </concept>
   <concept>
       <concept_id>10002951.10002952.10003190.10003195</concept_id>
       <concept_desc>Information systems~Query optimization</concept_desc>
       <concept_significance>500</concept_significance>
       </concept>
   <concept>
       <concept_id>10010147.10010178.10010179</concept_id>
       <concept_desc>Computing methodologies~Natural language processing</concept_desc>
       <concept_significance>300</concept_significance>
       </concept>
   <concept>
       <concept_id>10010147.10010178.10010219.10010220</concept_id>
       <concept_desc>Computing methodologies~Multi-agent systems</concept_desc>
       <concept_significance>100</concept_significance>
       </concept>
 </ccs2012>
\end{CCSXML}

\ccsdesc[500]{Information systems~Data management systems}
\ccsdesc[500]{Information systems~Query optimization}
\ccsdesc[300]{Computing methodologies~Natural language processing}
\ccsdesc[100]{Computing methodologies~Multi-agent systems}

\keywords{Agent memory, query optimization, large language models,
  self-tuning systems, cost-based optimization}

\maketitle

% ============ Sections ============
%!TEX root = ../main.tex
\section{Introduction}
\label{sec:intro}

%!TEX root = ../../main.tex
% Figure 1: motivating example (single column), drawn in TikZ.
\begin{figure}[t]
    \centering
    \begin{tikzpicture}[
        font=\footnotesize,
        >={Stealth[length=4.2pt,width=3.4pt]},
        qbox/.style={draw=darkgray_x11!50, rounded corners=2.5pt, fill=gray!7,
                     text width=2.25cm, align=left, inner sep=3.5pt,
                     execute at begin node={\hyphenpenalty=10000\exhyphenpenalty=10000}},
        plan/.style={draw=#1, line width=0.8pt, rounded corners=2.5pt, fill=#1!7,
                     text width=1.42cm, align=center, inner sep=3.5pt},
        op/.style={draw=#1, line width=0.8pt, rounded corners=2.5pt, fill=white,
                   text width=2.84cm, align=left, inner sep=3.5pt},
        band/.style={draw=darkgray_x11!50, rounded corners=2.5pt, fill=lightblue,
                     text width=7.92cm, align=center, inner sep=3.5pt},
        foot/.style={draw=darkgray_x11!50, rounded corners=2.5pt, fill=gray!10,
                     text width=7.92cm, align=center, inner sep=3.5pt},
        fb/.style={draw=cadmiumgreen, line width=0.8pt, rounded corners=2.5pt,
                   fill=lightgreen, text width=2.84cm, align=left, inner sep=3.5pt},
        flow/.style={->, draw=darkgray_x11!70, line width=0.8pt},
        loop/.style={->, draw=cadmiumgreen, line width=0.8pt, rounded corners=3pt,
                     dash pattern=on 2pt off 1.6pt},
    ]

    % ---- catalog band -------------------------------------------------
    \node[band] (cat) at (3.96,2.30)
      {\scriptsize Catalog $\mathrm{Cat}(\mathcal{H})$: entities, events,
       relations, attributes\,---\,indexes $\mathcal{I}_e,\mathcal{I}_a$};

    % ---- row 1: point lookup ------------------------------------------
    \node[qbox] (q1) at (1.125,1.00)
      {\textbf{Q1} \emph{``Who is my mortgage provider?''}};
    \node[plan=softblue] (p1) at (3.71,1.00)
      {\opPOINT\\[1pt]\scriptsize no quantifier};
    \node[op=softblue] (o1) at (6.50,1.00)
      {\scriptsize index scan $\mathcal{I}_a[\textsf{provider}]$\\
       $\Rightarrow$ \textbf{Raiffeisen Basel}\\[1pt]
       \scriptsize\textcolor{softblue}{1 LLM call}};

    % ---- second-order memory ------------------------------------------
    \node[fb] (fbk) at (6.50,-0.42)
      {\scriptsize \textbf{Second-order memory}\\
       log the hit entities, raise $\pi(e)$};

    % ---- row 2: aggregation -------------------------------------------
    \node[qbox] (q2) at (1.125,-1.85)
      {\textbf{Q2} \emph{``How many of my events involve outdoor activity?''}};
    \node[plan=softorange] (p2) at (3.71,-1.85)
      {\opAGG\\[1pt]\scriptsize ``how many''};
    \node[op=softorange] (o2) at (6.50,-1.85)
      {\scriptsize enumerate 8 preheated events on $\mathcal{I}_e$, tally yes/no\\
       $\Rightarrow$ \textbf{3}\\[1pt]
       \scriptsize two samples agree, third call skipped\\[1pt]
       \scriptsize\textcolor{softorange}{2 LLM calls}};

    % ---- footer -------------------------------------------------------
    \node[foot] (ft) at (3.96,-3.30)
      {\scriptsize \sys{}: \textbf{3} LLM calls in total\quad versus\quad
       Uniform-$K{=}3$: \textbf{6} calls};

    % ---- flow arrows ---------------------------------------------------
    \draw[flow] (q1.east) -- node[above,font=\scriptsize,inner sep=1.2pt]{$\rho$} (p1.west);
    \draw[flow] (p1.east) -- node[above,font=\scriptsize,inner sep=1.2pt]{$\Omega_c$} (o1.west);
    \draw[flow] (q2.east) -- node[above,font=\scriptsize,inner sep=1.2pt]{$\rho$} (p2.west);
    \draw[flow] (p2.east) -- node[above,font=\scriptsize,inner sep=1.2pt]{$\Omega_c$} (o2.west);
    \draw[flow] (cat.south -| o1.north) -- (o1.north);

    % ---- feedback loop: read path back to storage ----------------------
    \draw[loop] (o1.south) -- (fbk.north);
    \draw[loop] (fbk.south) -- (o2.north);
    \draw[loop] (fbk.west) -- (-0.26,-0.42) -- (-0.26,2.30) -- (cat.west);

    \end{tikzpicture}
    \caption{Motivating example. \sys{} selects a physical plan for each query
      (\opPOINT{} or \opAGG), serves it with the minimum LLM-call budget, and
      logs hit entities into second-order memory so that the next query reads
      them first. The interaction costs three LLM calls in total, where a
      Uniform-$K{=}3$ baseline would spend six.}
    \Description{Two queries from the same user flow left to right. A point
      lookup is dispatched to the point operator and answered in one LLM call;
      its hit entities are logged by the second-order memory, which raises
      their priority in the catalog so that the following aggregation query
      reads them first and is answered in two calls.}
    \label{fig:motivating}
\end{figure}


LLM agents now hold conversations that span weeks, and the interaction history of a single user grows to hundreds of thousands of tokens, in some deployments to millions~\cite{mem0,memgpt,zep}. Answering a question against that history is the core function of a personal agent, and it is also the function that current systems serve least well. The obstacle is not the cost of storing the history but the absence of any declared structure over it, since the history arrives as free-form natural language. A system therefore cannot tell a single-fact lookup apart from a count over scattered evidence, and cannot treat them differently. Every read becomes the same operation, a similarity scan over an opaque token stream followed by one generation call. The system never knows what a query is asking, what the same user has asked before, or what an answer ought to cost. Agent memory today is stored but not managed. Figure~\ref{fig:motivating} contrasts this behavior with \sys{} on a two-query interaction, and Figure~\ref{fig:scoreboard} previews the accuracy \sys{} reaches across eight long-term memory benchmarks.

\stitle{State-of-the-Art Approaches and Their Limits.}
Three lines of work attack this problem. Each improves how memory is written and leaves untouched how it is read, which is where a query is actually served. \emph{Long-context prompting} and \emph{RAG variants}~\cite{longcontext,rag} push raw history into the prompt or retrieve top-$k$ chunks by dense similarity. A point lookup and a cross-entity count then follow an identical retrieval-and-generation pipeline. \emph{Hierarchical memory systems} such as MemGPT~\cite{memgpt} page facts between context tiers, which decides \emph{where} a fact lives but not \emph{how} to compute over facts of different logical shapes. \emph{Structured memory systems} such as Mem0~\cite{mem0} and Zep~\cite{zep} invest heavily in write-time entity and graph extraction, yet both fall back to one dense-retrieval strategy at read time, and neither records which stored memories actually produced correct answers for a given user. The three lines differ in representation but share a common shape. Prior systems supply a storage layer and sometimes an access path, but no planner that distinguishes plan classes, no operator library that computes each class differently, and no cost model that prices an answer before paying for it.

The database community closed analogous gaps four decades ago by lifting flat files into \emph{managed data}, where schemas expose record shape, optimizers select physical plans over a declared algebra, and cost models allocate execution resources~\cite{selinger1979,graefe1993}. Recent systems carry that abstraction into LLM analytics over typed tables~\cite{lotus}, but the line stops where the data stops being typed. Agent memory is unstructured, its schema is latent in natural language, and its queries are user questions instead of SQL. No existing system transfers the optimization stack to memory itself.

\stitle{Challenges.}
Carrying the stack across that gap raises three obstacles. In each case a direct remedy exists and fails for a concrete reason.

\etitle{C1. Recovering a Plan Space without a Schema or Labels.}
Query optimization presupposes that the system can name the shape of a query, but agent memory offers neither a relational schema to plan against nor labeled examples to learn a planner from. Asking the LLM to tag each query would work, yet it spends an extra call on every question and turns the plan itself into a sampled variable, which is an undesirable property for a component every downstream decision depends on. Training a supervised planner is unavailable because no benchmark ships query-type labels.

\etitle{C2. Learning a User without Gold Supervision.}
A memory system should get better at a specific user as that user keeps asking, but nothing in the deployment reveals which stored records deserve priority. Loading every historical entity into the prompt exceeds the context window after a few sessions. Ranking by recency is available without any supervision, yet it ranks by the wrong quantity, because an entity a user returns to across months is old precisely because the relationship is long-standing. The signal that does matter, which records participate in correct answers, lies on the read path, and prior systems instrument only the write path.

\etitle{C3. Paying for Reasoning in Proportion to Difficulty.}
Replicated sampling buys accuracy on hard questions and wastes it on easy ones, yet difficulty is not observable before the answer exists. Lowering the budget uniformly sacrifices the queries that needed it, and an LLM-based difficulty predictor reintroduces the per-query call the budget was meant to save.

\stitle{Our Proposal.}
We present \sys{}, a memory management system that answers all three obstacles with one move, by giving agent memory the storage, planning, and execution separation of a DBMS. Three mechanisms follow from that separation.

\etitle{A Typed Catalog with Zero-Supervision Plan Selection.}
\sys{} lifts sessions into a catalog of entities, events, relations, and attributes, materialized lazily on first access, and routes each question to one of three physical plans using only surface form, the quantifier and comparative markers in the question and the shape of the answer choices. The planner costs no LLM call and needs no labels, which resolves C1. Three operators then execute the chosen plan over the catalog indexes. A count enumerates and tallies where a comparison extracts and sorts, instead of both being handed to one generic prompt.

\etitle{Second-Order Memory over the Read Path.}
\sys{} records, per user, which plan classes have been issued and which catalog entities appeared in accepted answers, then folds that record back into index priority and grounding-block composition. The system observes the read path that prior work discards, which resolves C2, and it transfers workload-driven index recommendation~\cite{chaudhuri1997autotuning} from self-tuning databases to agent memory.

\etitle{Agreement as a Cost Signal.}
\sys{} assigns a call budget per plan class and then truncates it adaptively by comparing two independent samples. Reading sample agreement as a stopping signal is known~\cite{adaptiveconsistency,escsc}, but it is applied to one query at a time and therefore cannot avoid the first two samples. \sys{} determines from the plan class alone that a point lookup will not benefit from a second sample, and it therefore allocates across classes before adapting within one, which resolves C3.

\begin{example}
\label{exam:intro}
Alice asks her agent \emph{``How many of my events involve outdoor activity?''} after six earlier questions in the same session. A Full-Context baseline scans the whole history and answers 5, miscounting evidence scattered across sessions, where the correct answer is 3. A graph memory retrieves six outdoor-related nodes but never resolves the outdoor predicate, and the model guesses 4. \sys{} routes the question to its aggregation plan, surfaces the event entities that the six earlier questions already marked as relevant, enumerates each candidate with an explicit yes-or-no decision, and answers 3. The two samples it draws agree, and it stops after two calls. Each of the three failure modes above is removed by a different one of the three mechanisms. Table~\ref{tbl:casestudy} gives the full trace.
\end{example}

%!TEX root = ../../main.tex
% Radar chart with per-axis min-max normalization: on each axis,
% Full-Context maps to r=0.15R and MemCatalog maps to r=1.00R (regular octagon),
% Prior Best falls linearly in between at r = 0.15 + 0.85*(prior-fc)/(sys-fc).
% Raw accuracy is annotated at every vertex.
\begin{figure}[t]
    \centering
    \begin{tikzpicture}[every node/.style={font=\footnotesize}]
        \filldraw[fill=gray!35, fill opacity=0.4, draw=gray!80, thick] (0,0) rectangle (0.35,0.18);
        \node[anchor=west] at (0.4,0.09) {Full-Context};
        \filldraw[fill=softblue, fill opacity=0.28, draw=softblue, very thick] (2.6,0) rectangle (2.95,0.18);
        \node[anchor=west] at (3.0,0.09) {Prior Best};
        \filldraw[fill=dropred, fill opacity=0.32, draw=dropred, very thick] (5.0,0) rectangle (5.35,0.18);
        \node[anchor=west] at (5.4,0.09) {\textbf{\sys{} (ours)}};
    \end{tikzpicture}

    \resizebox{0.95\columnwidth}{!}{%
    \begin{tikzpicture}[
        every node/.style={font=\footnotesize},
        axis/.style={gray!45, thin},
        grid/.style={gray!25, dashed, very thin},
        v/.style={circle, inner sep=0pt, minimum size=3.4pt},
        vfc/.style={v, fill=gray!85},
        vpr/.style={v, fill=softblue},
        vmc/.style={v, fill=dropred},
        lblfc/.style={font=\tiny, gray!75, inner sep=0.5pt},
        lblpr/.style={font=\tiny, softblue, inner sep=0.5pt},
        lblmc/.style={font=\tiny\bfseries, dropred, inner sep=0.5pt},
    ]
        \def\R{3.4}
        \foreach \a in {90,45,0,-45,-90,-135,180,135}
            \draw[axis] (0,0) -- (\a:\R);
        \foreach \r in {0.33, 0.66, 1.0}
            \draw[grid] (90:{\r*\R}) -- (45:{\r*\R}) -- (0:{\r*\R}) -- (-45:{\r*\R})
                       -- (-90:{\r*\R}) -- (-135:{\r*\R}) -- (180:{\r*\R}) -- (135:{\r*\R}) -- cycle;
        \node[above=0.10cm]        at (90:\R)   {\textbf{\locomo}};
        \node[above right=0.03cm]  at (45:\R)   {\textbf{\lmev}};
        \node[right=0.10cm]        at (0:\R)    {\textbf{\membench}};
        \node[below right=0.03cm]  at (-45:\R)  {\textbf{\lmeo}};
        \node[below=0.10cm]        at (-90:\R)  {\textbf{\lmes}};
        \node[below left=0.03cm]   at (-135:\R) {\textbf{\lmem}};
        \node[left=0.10cm]         at (180:\R)  {\textbf{\evomem}};
        \node[above left=0.03cm]   at (135:\R)  {\textbf{\dynamicmem}};
        %================ Full-Context (gray): innermost octagon at 0.15R =================%
        \filldraw[fill=gray!35, fill opacity=0.32, draw=gray!80, thick, line join=round]
            (90:{0.15*\R})    -- (45:{0.15*\R})   -- (0:{0.15*\R})    -- (-45:{0.15*\R})
         -- (-90:{0.15*\R})   -- (-135:{0.15*\R}) -- (180:{0.15*\R})  -- (135:{0.15*\R})
         -- cycle;
        %================ Prior Best (blue): per-axis interpolation =================%
        % ratios (r_prior): LoCoMo 0.922, LMEv 0.868, MB 0.777, LMEo 0.763, LMEs 0.777, LMEm 0.751, EvoMem 0.731, DynMem 1.000
        \filldraw[fill=softblue, fill opacity=0.24, draw=softblue, very thick, line join=round]
            (90:{0.922*\R})   -- (45:{0.868*\R})  -- (0:{0.777*\R})   -- (-45:{0.763*\R})
         -- (-90:{0.777*\R})  -- (-135:{0.751*\R}) -- (180:{0.731*\R}) -- (135:{1.000*\R})
         -- cycle;
        %================ MemCatalog (red): outermost regular octagon =================%
        \filldraw[fill=dropred, fill opacity=0.28, draw=dropred, very thick, line join=round]
            (90:\R)   -- (45:\R)  -- (0:\R)   -- (-45:\R)
         -- (-90:\R)  -- (-135:\R) -- (180:\R) -- (135:\R)
         -- cycle;
        %================ Vertex markers + raw accuracy labels =================%
        \node[vfc] at (90:{0.15*\R}) {}; \node[lblfc, right] at (90:{0.15*\R}) {\,52.3};
        \node[vpr] at (90:{0.922*\R}) {}; \node[lblpr, right] at (90:{0.922*\R}) {\,58.2};
        \node[vmc] at (90:\R) {};         \node[lblmc, left]  at (90:\R) {58.8\,};
        \node[vfc] at (45:{0.15*\R}) {}; \node[lblfc, below right] at (45:{0.15*\R}) {4.5};
        \node[vpr] at (45:{0.868*\R}) {}; \node[lblpr, below right] at (45:{0.868*\R}) {9.4};
        \node[vmc] at (45:\R) {};         \node[lblmc, above left]  at (45:\R) {10.3};
        \node[vfc] at (0:{0.15*\R}) {};  \node[lblfc, above] at (0:{0.15*\R}) {86.7};
        \node[vpr] at (0:{0.777*\R}) {};  \node[lblpr, above] at (0:{0.777*\R}) {91.2};
        \node[vmc] at (0:\R) {};          \node[lblmc, below] at (0:\R) {92.8};
        \node[vfc] at (-45:{0.15*\R}) {}; \node[lblfc, above right] at (-45:{0.15*\R}) {49.3};
        \node[vpr] at (-45:{0.763*\R}) {}; \node[lblpr, above right] at (-45:{0.763*\R}) {62.2};
        \node[vmc] at (-45:\R) {};         \node[lblmc, below left]  at (-45:\R) {67.2};
        \node[vfc] at (-90:{0.15*\R}) {}; \node[lblfc, right] at (-90:{0.15*\R}) {\,43.6};
        \node[vpr] at (-90:{0.777*\R}) {}; \node[lblpr, right] at (-90:{0.777*\R}) {\,59.3};
        \node[vmc] at (-90:\R) {};         \node[lblmc, left]  at (-90:\R) {64.9\,};
        \node[vfc] at (-135:{0.15*\R}) {}; \node[lblfc, above left] at (-135:{0.15*\R}) {58.5};
        \node[vpr] at (-135:{0.751*\R}) {}; \node[lblpr, above left] at (-135:{0.751*\R}) {70.3};
        \node[vmc] at (-135:\R) {};         \node[lblmc, below right] at (-135:\R) {75.2};
        \node[vfc] at (180:{0.15*\R}) {}; \node[lblfc, below] at (180:{0.15*\R}) {16.7};
        \node[vpr] at (180:{0.731*\R}) {}; \node[lblpr, below] at (180:{0.731*\R}) {22.1};
        \node[vmc] at (180:\R) {};         \node[lblmc, above] at (180:\R) {24.6};
        \node[vfc] at (135:{0.15*\R}) {}; \node[lblfc, below left] at (135:{0.15*\R}) {33.5};
        \node[vpr] at (135:{1.000*\R}) {}; \node[lblpr, below left] at (135:{1.000*\R}) {34.9};
        \node[vmc] at (135:\R) {};         \node[lblmc, above right] at (135:\R) {34.9};
    \end{tikzpicture}
    }
    \caption{Per-benchmark accuracy on 8 long-term memory benchmarks. Each axis uses its own linear scale: Full-Context maps to the innermost ring, \sys{} to the outermost regular octagon, and Prior Best falls at its position between them, so that within-benchmark differences among the three systems are directly visible. Raw accuracy (\%) is annotated at every vertex. The prior-best system varies per benchmark (\membench: Mem0; \lmev: LongMemEval-Fw; the remaining six: Zep), matching Table~\ref{tbl:main_results}.}
    \Description{Radar chart over eight long-term memory benchmarks. On each
      axis MemCatalog lies on the outermost octagon, the strongest prior memory
      system lies between it and the Full-Context baseline on the innermost
      ring, and the raw accuracy of all three systems is annotated at each
      vertex.}
    \label{fig:scoreboard}
\end{figure}


\stitle{Contributions.} (1) We formulate agent memory management as query optimization and build \sys{}, which supplies the catalog, planner, and operator library that prior memory systems lack (Sections~\ref{sec:overview} and~\ref{sec:framework}).
(2) We introduce second-order memory, which closes a feedback loop from the read path back to physical storage and lets a memory system improve on a user without supervision (Section~\ref{sec:secondorder}).
(3) We combine class-level budget allocation with agreement-based stopping, which together reach a cost point that neither mechanism attains alone (Section~\ref{sec:executor}).
(4) We evaluate \sys{} on eight long-term memory benchmarks and eight LLM backbones, where it matches or improves the strongest prior memory system on every benchmark while using fewer LLM calls than a uniform sampling reference of equal plan depth, and we release code, harness, and per-entry outputs for full reproduction\footnote{Materials will be made publicly available upon acceptance.}.

%!TEX root = ../main.tex
\section{Preliminaries}
\label{sec:preliminary}

This section fixes notation, defines the agent memory management problem, describes how any memory-augmented LLM agent processes a query, and lists the three access primitives that structure the rest of the paper.

\subsection{Interaction History and Memory Query}
\label{subsec:prelim-data}

\stitle{Interaction History.}
An LLM agent interacts with a user across a sequence of \emph{sessions} $\mathcal{H}=\{s_1,\dots,s_N\}$, where $s_i=(t_i,\mathrm{txt}_i)$ is a timestamp $t_i$ paired with a natural-language transcript $\mathrm{txt}_i$ of mixed user and assistant turns. Sessions are read-only after they close, and $\mathcal{H}$ is unbounded and grows monotonically over time. In our benchmarks a single user's history reaches up to $1.5$M tokens spread over tens of sessions~\cite{longmemeval,locomo}, well beyond any current context window at inference time.

\stitle{Memory Query.}
A memory query is a triple $q=(\mathrm{Q},\mathcal{C},\tau)$, where $\mathrm{Q}$ is a natural-language question, $\mathcal{C}$ is an optional set of answer choices (empty for open-ended questions), and $\tau$ is the timestamp at which the query is issued. The agent must return an answer $\hat{a}$ derivable from the history available at that moment, $\mathcal{H}_{\leq\tau}=\{s_i:t_i\leq\tau\}$. Answers are scored against a gold reference, with accuracy or F1 computed on $\hat{a}$ (Section~\ref{subsec:setup}).

\subsection{Problem Statement}
\label{subsec:prelim-problem}

Given a stream of queries $q_1,q_2,\dots$ over an evolving $\mathcal{H}$, an \emph{agent memory management system} produces answers $\hat{a}_1,\hat{a}_2,\dots$ under three simultaneous objectives.
\begin{enumerate}
\item[\textbf{(O1)}] \textbf{Accuracy.} Maximize the fraction of $\hat{a}_t$ that match the gold answer for $q_t$.
\item[\textbf{(O2)}] \textbf{Bounded per-query cost.} Keep the per-query LLM cost $\mathrm{cost}(q_t)$ bounded and manageable as $|\mathcal{H}|$ grows, letting the system scale to long histories and heavy usage.
\item[\textbf{(O3)}] \textbf{Improvement over time.} Later queries on the same $\mathcal{H}$ should improve, or at least not degrade, given the information gathered from earlier ones.
\end{enumerate}
Prior memory systems have focused overwhelmingly on \textbf{(O1)}~\cite{mem0,zep,memgpt,longmemeval}: their headline results report accuracy on a single test query per user, with cost and per-user adaptation treated implicitly or not at all. \sys{} treats \textbf{(O2)} and \textbf{(O3)} as first-class objectives on par with \textbf{(O1)}, and enforces each by a dedicated architectural layer, the executor for cost (Section~\ref{sec:executor}) and the second-order memory for per-user improvement (Section~\ref{sec:secondorder}).

\subsection{Memory Operations in LLM Agents}
\label{subsec:prelim-flow}

Regardless of the particular memory system, an LLM agent processes each user interaction along two coupled paths. A \emph{write path} updates the memory store when a new session closes, and a \emph{read path} answers a query using the accumulated store. Both paths spend LLM calls, and both are what \sys{} instruments.

Every memory system, whatever its internal representation, runs the same two paths. A \emph{write path} fires when a session closes, converting $s_i$ into the store's representation and appending it to a memory store $\mathcal{M}$. Systems differ here in representation alone, with long-context prompting keeping $s_i$ verbatim, RAG embedding chunks of it, Mem0 extracting facts into a dictionary, and Zep building temporal-graph edges. A \emph{read path} fires per query and proceeds in three steps that are common to all of them, \textbf{(R1)} retrieving a candidate set from $\mathcal{M}$, \textbf{(R2)} packing candidates into a grounding block that fits the context window, and \textbf{(R3)} reasoning over that block to produce $\hat{a}_t$.

Two observations about this structure motivate the rest of the paper. First, $\mathrm{cost}(q_t)$ is dominated by R3, and replicated sampling multiplies it on every query without regard to whether the query needed the samples. Second, prior systems run an identical R1--R3 pipeline for every query and retain nothing from the outcome. Neither the shape of $\mathrm{Q}_t$ nor the history of this user's earlier queries changes how the next one is served.

\subsection{Memory Access Primitives}
\label{subsec:prelim-primitives}

Empirically, natural-language queries against agent memory fall into three access patterns that differ in the shape of the computation they require. This taxonomy is the entry point for both the planner and the physical operator library of Section~\ref{sec:framework}.

\begin{itemize}[leftmargin=*]
\item \textbf{Point lookup (\opPOINT).} Retrieve a single stored fact by attribute key, \eg \emph{``Who is my mortgage provider?''}. The answer is a single value from the catalog's attribute index and does not require aggregation or comparison.
\item \textbf{Aggregation (\opAGG).} Compute a quantity (typically a count) over multiple records, \eg \emph{``How many of my events involve outdoor activity?''}. The answer requires enumerating candidates from the entity index and applying a predicate to each.
\item \textbf{Comparison (\opCMP).} Rank or order two or more records on a shared attribute, \eg \emph{``Which lasted longer, the Basel trip or the Zurich conference?''}. The answer requires extracting the compared key from each candidate and reducing it to an $\argmax$.
\end{itemize}

These three classes have direct DBMS analogues (\emph{index scan}, \emph{hash-aggregate}, \emph{sort-merge}) that Sections~\ref{subsec:planner} and~\ref{subsec:budget} exploit, and the LLM's failure modes differ across them (reliable point lookup, undercount on scattered aggregation~\cite{ruler}, and misalignment on cross-passage comparison). Prior memory systems collapse the three into one retrieve-then-generate pipeline and inherit all three failure modes. \sys{} splits them and gives each its own operator.

%!TEX root = ../main.tex
\section{System Overview}
\label{sec:overview}

%!TEX root = ../../main.tex
% Placeholder: framework overview figure (two-column span), to be hand-drawn.
\begin{figure*}[t]
    \centering
    \setlength{\fboxrule}{0.8pt}%
    \textcolor{gray!55}{\framebox[\dimexpr\textwidth-2\fboxsep-2\fboxrule\relax]{%
      \begin{minipage}[c][0.24\textwidth][c]{0.95\textwidth}
        \centering
        \textcolor{gray!65}{\normalsize\itshape [Placeholder] Framework overview of \sys{}}\\[0.6em]
        \textcolor{gray!55}{\small catalog $+$ planner $+$ operators
          $\;\longrightarrow\;$ second-order memory $\;\longrightarrow\;$ cost-aware executor}
      \end{minipage}}}
    \caption{Overview of \sys{}. A foundational framework (catalog, planner,
      operators) is augmented by two optimization layers:
      the Second-Order Memory, which closes a workload feedback loop,
      and the Cost-Aware Executor, which combines pattern-aware budgeting with
      agreement-based early stopping.}
    \Description{Placeholder box for the system architecture diagram of
      MemCatalog, showing the catalog, planner, and operators
      as the base framework, with the second-order memory and the cost-aware
      executor layered on top.}
    \label{fig:framework_overview}
\end{figure*}


This section traces one query end to end, so that the division of labor among the three mechanisms is visible before any of them is specified in detail. Figure~\ref{fig:framework_overview} shows the resulting architecture, in which a framework supplies the plan and two optimization layers refine what that plan reads and what it costs.

\stitle{Life Cycle of a Query.}
Given a query $q$, \sys{} evaluates Eq.~(\ref{eq:sys}) in four steps. It first \emph{plans} the query, mapping $q$ to one of \point, \agg, or \cmp{} from surface form alone (Section~\ref{subsec:planner}). It then \emph{grounds} the selected plan, resolving the catalog indexes $\mathcal{I}_e$ and $\mathcal{I}_a$ that the chosen operator will read, with the second-order memory reweighting those indexes by the per-user priority $\pi(e)$ so that entities this user has already confirmed appear first (Section~\ref{subsec:feedback}). It then \emph{executes}, issuing LLM calls under the operator's fixed template while the executor decides how many calls the plan class warrants and stops early once two samples agree (Sections~\ref{subsec:budget} and~\ref{subsec:shortcircuit}). It finally \emph{records} the outcome, writing the entities that appeared in the accepted answer into the workload log $\mathcal{W}_u$ so that the next query starts from a better-ordered catalog (Section~\ref{subsec:log}).

\stitle{Interlocking of the Three Mechanisms.}
The ordering above is not a pipeline of independent stages, since each mechanism consumes something another produces. The executor cannot price a query until the planner has named its class. The planner cannot read a useful candidate set until the second-order memory has ranked the catalog. The second-order memory cannot update until the executor has committed to an answer worth recording. Removing any one therefore degrades the other two, which Section~\ref{sec:experiments} confirms by ablation. The remainder of the paper takes the three in turn.

%!TEX root = ../main.tex
\section{Memory as Managed Data}
\label{sec:framework}

This section describes how \sys{} turns unstructured agent memory into managed data. Prior systems process every query through one generic retrieval pipeline, regardless of what the query asks. \sys{} factors the computation into two parts: a \emph{catalog} that gives the interaction history a structure (Section~\ref{subsec:catalog}), and a \emph{planner} that routes each query to one of three specialized operators (Section~\ref{subsec:planner}). The overall computation is:
\begin{equation}
\label{eq:sys}
\hat{a}_t \;=\; \underbrace{\Omega_{c}}_{\text{Operator}}\bigl(q_t,\ \underbrace{\mathrm{Cat}(\mathcal{H})}_{\text{Catalog}}\bigr),
\qquad c = \underbrace{\rho(q_t)}_{\text{Planner}},
\end{equation}
where $\rho$ selects a query pattern $c\in\{\point,\agg,\cmp\}$ and $\Omega_c$ is the operator that handles pattern $c$. We continue to use Alice from Example~\ref{exam:intro} as the running illustration.

\subsection{Typed Catalog}
\label{subsec:catalog}

To specialize query processing, the system first needs to know the shape of its data. Raw conversation transcripts do not reveal any shape. The catalog recovers that shape by extracting typed records from sessions.

\stitle{Typed Records.}
The catalog $\mathrm{Cat}(\mathcal{H})$ decomposes each session $s_i$ into four record types: \emph{entities} $\mathcal{E}$ (named actors, objects, locations), \emph{events} $\mathcal{V}$ (timestamped occurrences), \emph{relations} $\mathcal{R}\subseteq\mathcal{E}\times P\times\mathcal{E}$ (entity pairs with named predicates), and \emph{attributes} $\mathcal{A}\subseteq\mathcal{E}\times K\times\mathrm{Val}$ ((entity, key, value) triples). For example, a session in which Alice mentions opening a mortgage contributes the entities $\{\textsf{Alice},\textsf{Raiffeisen Basel}\}$, the event $\textsf{(opened mortgage, 2025-03-14)}$, the relation $\textsf{(Alice, has-provider, Raiffeisen Basel)}$, and the attribute $\textsf{(mortgage, rate, 2.4\%)}$. What was opaque text now has a visible structure.

\stitle{On-Demand Extraction.}
Records are extracted only when needed. When a query first touches session $s_i$, one templated LLM call returns $(\mathcal{E}_i,\mathcal{V}_i,\mathcal{R}_i,\mathcal{A}_i)$, and the result is cached. Extraction runs once per session and is amortized over all later queries that touch the same session, in the same way a DBMS builds an index once and reuses it across a workload.

\stitle{Access Paths.}
Operators need two kinds of lookup: finding which sessions mention a given entity, and finding the value of a given attribute. The catalog exposes one index for each: an entity index $\mathcal{I}_e$ that maps each entity to the sessions it appears in, and an attribute index $\mathcal{I}_a$ that maps each attribute key to its (entity, value, session) records. Operators read memory only through these two indexes, which keeps the storage swappable and the ranking tunable by the second-order memory (Section~\ref{sec:secondorder}).

\subsection{Query Classification and Operators}
\label{subsec:planner}

With structured data in place, the next question is which computation a given query needs. This subsection answers it in two steps: a classification rule that names the query pattern, and an operator library that executes each pattern differently.

\stitle{Query Classification.}
The planner $\rho$ uses only the query text and the answer choices:
\begin{equation}
\label{eq:classifier}
\rho(q) \;=\; \begin{cases}
\agg, & q \text{ contains a phrase in } \mathrm{QuantPhr}\ \lor\ \mathcal{C} \text{ are numeric}, \\
\cmp, & \text{else if } q \text{ contains a phrase in } \mathrm{CmpPhr}, \\
\point, & \text{otherwise,}
\end{cases}
\end{equation}
where $\mathrm{QuantPhr}$ is a fixed set of quantifier phrases (\textsf{how many}, \textsf{the number of}, \textsf{count}) and $\mathrm{CmpPhr}$ is a fixed set of comparative markers (\textsf{more}, \textsf{longer}, \textsf{most}). No LLM call or labeled data is needed. For example, \emph{``Who is my mortgage provider?''} falls through to \point{}. \emph{``How many of my events involve outdoor activity?''} matches \textsf{how many} and routes to \agg{}. \emph{``Which lasted longer, the Basel trip or the Zurich conference?''} matches \textsf{longer} and routes to \cmp{}.

\stitle{Why Misclassification Is Safe.}
A misclassified query affects only accuracy, not correctness. All three operators read the same catalog through the same indexes, so every operator retrieves admissible evidence and returns a well-formed answer regardless of which one is selected. Moreover, \point{} is the fall-through default, so unrecognized queries take the cheapest path. Section~\ref{subsec:exp7routing} measures the accuracy that classification contributes.

\stitle{Operator Library.}
A point query needs a single value from the catalog, an aggregation query needs to enumerate candidates and count matches, and a comparison query needs to retrieve one value per candidate and compare them. Each operator implements that computation with the catalog indexes:
\begin{align*}
\Omega_{\point}(q) &\;=\; \mathrm{LLM}\bigl(\mathrm{template}_{\point}(\mathcal{I}_a[k_q])\bigr), \\
\Omega_{\agg}(q)   &\;=\; \bigl|\{\,e\in \mathcal{I}_e[T_q]\ :\ \mathrm{LLM}(\mathrm{pred}_q,e)=\texttt{yes}\}\bigr|, \\
\Omega_{\cmp}(q)   &\;=\; \argmax_{e\in\{e_1,e_2\}}\ \mathcal{I}_a[k_q](e).
\end{align*}
\opPOINT{} looks up the requested key $k_q$ in $\mathcal{I}_a$ and returns the value directly. \opAGG{} enumerates candidate entities of the referenced type $T_q$ from $\mathcal{I}_e$, evaluates each against the query predicate with a bounded \texttt{yes}/\texttt{no} judgment, and emits an integer count. On the outdoor query, it lists Alice's eight events, tags each, and emits \texttt{Answer: 3}. This makes the count explicit, directly enumerable from structured evidence. \opCMP{} extracts the compared key $k_q$ (here \textsf{duration}) for both candidates from $\mathcal{I}_a$ and presents them side by side, turning an alignment problem over long contexts into a direct comparison of two values.

\stitle{Remarks.}
The catalog, the planner, and the operator library form the foundational framework of \sys{}. Making the query pattern $c=\rho(q)$ explicit is the precondition for the two optimization layers that follow: the second-order memory tunes the indexes that operators read (Section~\ref{sec:secondorder}), and the cost-aware executor sets a different budget for each query pattern (Section~\ref{sec:executor}). Neither layer would be possible without knowing the query pattern first.

%!TEX root = ../main.tex
\section{Second-Order Memory over the Read Path}
\label{sec:secondorder}

The framework of Section~\ref{sec:framework} evaluates each query in isolation. The hundredth query for a user is therefore served in the same way as the first. Prior memory systems share this limitation. They observe only the \emph{write path}, which sessions arrived, and never the \emph{read path}, which stored records answered which queries. \sys{} closes this loop with a per-user workload log $\mathcal{W}_u$ (Section~\ref{subsec:log}) and a feedback rule that reshapes the catalog's physical organization (Section~\ref{subsec:feedback}).

\subsection{Recording the Read Path}
\label{subsec:log}

Adapting storage to a user presupposes a record of what that user's queries actually consumed. Prior systems never keep that record. This subsection specifies what \sys{} logs, and it keeps the log small enough that logging does not itself become a scaling problem. For each user $u$, \sys{} maintains a compact log of query outcomes
\[
\mathcal{W}_u \;=\; \{\,(q_t,\ \rho(q_t),\ E_t^{+},\ \tau_t)\,\}_{t\leq T},
\]
where $E_t^{+}\subseteq\mathcal{E}$ is the set of catalog entities that appeared in the accepted answer $\hat{a}_t$ or its supporting evidence. The log carries three signals. The sequence of plan classes $\{\rho(q_t)\}$ describes \emph{what the user asks}. The sets $\{E_t^{+}\}$ record \emph{what memory answers it}. Sessions left unread beyond a staleness horizon indicate \emph{what memory is dead weight}. After Alice's first six queries, her log records a \point-heavy class mix over financial and scheduling topics, hits on $\textsf{Raiffeisen Basel}$ and her calendar events, and no reads on her older travel sessions.

\stitle{Per-Entity Statistics.}
The log is statistical, not verbatim. For each catalog entity $e$ we keep three running numbers only, an issue count $n(e)=|\{t\leq T: e\in q_t\}|$ (how often $e$ has been mentioned by any query), an answer-hit count $h(e)=|\{t\leq T: e\in E_t^{+}\}|$ (how often $e$ appeared in an accepted answer), and a last-hit time $\tau^{\star}(e)=\max\{\tau_t: e\in E_t^{+}\}$. Storage is therefore $O(|\mathcal{E}|)$ per user and independent of query volume $T$. A DBMS keeps histograms and access counters instead of a replay log of every statement, for the same reason, since the optimizer needs the workload's distribution and not its transcript.

\stitle{What Counts as an Accepted Answer.}
No gold labels exist at deployment time, and acceptance must therefore be defined operationally. \sys{} marks an answer accepted when the executor commits to it under Eq.~(\ref{eq:short}), which happens when two independent samples agree, and records only the entities the operator read in producing it. The score moves a ranking and never a decision, and a promoted entity changes which records the operator sees first, never what it computes or may return. The loop is a retrieval prior, not a cache of asserted facts.

\stitle{Per-User Isolation.}
$\mathcal{W}_u$ and every quantity derived from it are strictly per-user, with no cross-user aggregation. A new user starts from a neutral prior and adapts as queries arrive.

\subsection{Reshaping Storage from Observed Workload}
\label{subsec:feedback}

A log is useful only if something reads it. This subsection turns the log into a single per-entity score and then spends that score on three decisions about where memory physically sits, leaving query semantics untouched throughout.

\stitle{Workload-Driven Priority Score.}
Each catalog entity $e$ carries a priority score
\begin{equation}
\label{eq:pi}
\pi(e) \;=\; \underbrace{\alpha\cdot h_r(e)}_{\text{used in answers}}\;+\;\underbrace{\beta\cdot r(e)}_{\text{used recently}}\;+\;\underbrace{\gamma\cdot c(e)}_{\text{matches the class}},
\end{equation}
where $h_r(e)=h(e)/n(e)$ is the hit rate on queries that mention $e$, $r(e)=\exp\bigl(-\lambda(\tau_{\mathrm{now}}-\tau^{\star}(e))\bigr)$ is a recency term with decay rate $\lambda$, and $c(e)$ is the fraction of queries mentioning $e$ whose plan class matches $e$'s record type. The three weights are fixed at $(\alpha,\beta,\gamma)=(0.5,0.3,0.2)$ for every benchmark and backbone we report, with direct workload evidence given the dominant share and the remaining two terms acting as tie-breakers. We do not tune them per dataset, and Section~\ref{subsec:exp5} reports the effect of removing each term outright.

\stitle{Storage-Level Actions.}
$\pi(e)$ is refreshed after every query and drives the physical organization of the catalog through three actions, governed by fixed system-wide thresholds. \emph{Index warming} keeps the highest-priority entities by $\pi$ in a resident cache that biases candidate retrieval. \emph{Grounding-block composition} assembles the evidence block prepended to an operator's prompt from the top candidate entities by $\pi(e)$, and not from a static set. \emph{Session eviction} compresses or drops any session $s_i$ whose entities all fall below a cold threshold $\pi_{\mathrm{cold}}$. By Alice's seventh query, her calendar events sit at the top of the grounding block while the never-read travel sessions have been compressed away. When \opAGG{} then enumerates candidates for the outdoor question the eight relevant events are the first records it reads. The operator itself is unchanged, and the workload has reorganized only the catalog it reads.

All three actions operate on \emph{where} memory sits, not on \emph{what} a query means. The access paths $\mathcal{I}_e,\mathcal{I}_a$ simply return workload-reweighted results, and every physical operator benefits without issuing an additional LLM call.

\stitle{Remarks.}
The priority score of Eq.~(\ref{eq:pi}) implements the self-tuning-index principle~\cite{chaudhuri1997autotuning,chaudhuri2007selftuning} and the workload-adaptive stance of learned optimizers~\cite{neoopt,baoopt,balsa}, transferred to agent memory, where the system observes the query log, updates statistics online, reshapes physical storage, and leaves query semantics untouched. In \sys{} the recommender is $\pi(e)$, the log is $\mathcal{W}_u$, and the storage is the catalog. To our knowledge, no prior agent memory system closes this workload-to-storage loop. This is why prior accuracy on a given user stays flat as that user issues more queries (Section~\ref{subsec:exp6}).

%!TEX root = ../main.tex
\section{Cost-Aware Execution}
\label{sec:executor}

A memory system that answers well but spends unboundedly does not scale to long histories and heavy usage. This section prices what \sys{} spends and then shows how to spend less of it without giving up accuracy.

\stitle{Accounting.}
\sys{} issues LLM calls on both the write path (catalog extraction) and the read path (query reasoning). Write-path cost is paid once per session and amortizes across queries. The rest of this section concerns read-path cost, which dominates. The standard remedy for accuracy, self-consistency at $K{=}3$~\cite{sc}, triples that cost on \emph{every} query, regardless of whether the extra samples change the answer. Formally, the total cost over a workload of $Q$ queries with class fractions $(p_{\point},p_{\agg},p_{\cmp})$ under fixed budgets $(K_{\point},K_{\agg},K_{\cmp})$ is
\begin{equation}
\label{eq:cost}
C \;=\; Q\cdot\bigl(p_{\point}K_{\point} + p_{\agg}K_{\agg} + p_{\cmp}K_{\cmp}\bigr).
\end{equation}
A Uniform-$K{=}3$ reference sets $K_{\point}=K_{\agg}=K_{\cmp}=3$, so $C=3Q$.

The executor makes this budget an explicit, plan-aware decision, exploiting two independent sources of waste. One is the class of queries saturated at a single call (Section~\ref{subsec:budget}). The other is the individual query whose answer is stable before its budget is spent (Section~\ref{subsec:shortcircuit}). The two act on disjoint slices, and their savings compose.

\subsection{Plan-Class Budgeting}
\label{subsec:budget}

The first source of waste is structural, since one class of query never benefits from the samples it is given. This subsection assigns a budget per class and derives what that assignment saves as a function of workload composition. The executor pre-declares a per-class budget instead of a per-query one, setting $K_{\point}=1$ and $K_{\agg}=K_{\cmp}=3$. The asymmetry follows from what each plan computes, not from a heuristic. A \point{} query resolves a single stored fact. Once $\textsf{(Alice, has-provider, Raiffeisen Basel)}$ is in the grounding block, replicated sampling can only re-read it, and a single call is both necessary and sufficient. \agg{} and \cmp{} queries compose evidence across multiple records, as when the outdoor count must decide eight independent \texttt{yes}/\texttt{no} judgments, and there sampling genuinely reduces variance, so these classes keep the full budget. Substituting these budgets into Eq.~(\ref{eq:cost}), the saving against Uniform-$K{=}3$ is $\Delta C=2p_{\point}Q$. On \membench{} with $p_{\point}=0.61$, this alone reduces $C$ from $900$ to $532$, a $40.9\%$ cut. The saving grows with the share of point queries and costs nothing in accuracy on exactly the queries it economizes, in line with the DBMS discipline of pricing an operator by what it must compute and not by a uniform per-tuple rate.

\subsection{Stability Short-Circuit}
\label{subsec:shortcircuit}

The second source of waste lies inside the hard classes, where some instances are easy but indistinguishable from the rest before execution. This subsection recovers that distinction at no extra cost by reading the agreement between samples the system was going to draw anyway. Within the \agg/\cmp{} budget, the executor decides adaptively whether the third call is worth making. It draws two independent samples $p_1,p_2$ under the operator's template, measures their token-level agreement $\mathrm{agree}(\cdot,\cdot)\in[0,1]$ against a threshold $\theta$, and draws a third sample $p_3$ only when the two disagree, giving the rule
\begin{equation}
\label{eq:short}
\hat{a} \;=\; \begin{cases}
p_1, & \mathrm{agree}(p_1,p_2)\geq\theta, \\
\mathrm{Majority}(p_1,p_2,p_3), & \text{otherwise.}
\end{cases}
\end{equation}
Under Eq.~(\ref{eq:short}) on Alice's outdoor question, both samples enumerate the same eight events and emit \texttt{Answer: 3}, so the agreement is high, the third call is skipped, and the query costs two. Had one sample returned \texttt{3} and the other \texttt{4}, the low agreement would have triggered escalation, spending the third call where it can change the outcome. We fix $\theta=0.6$ for all benchmarks and all backbones, and it is never retuned per dataset. Combined with plan-class budgeting, agreement-based stopping drives the total cost on \membench{} from $532$ to $444$, a further $16.5\%$ cut and $50.7\%$ overall against the Uniform-$K{=}3$ reference.

\stitle{Agreement as a Cost Signal.}
The reinterpretation here is the contribution. Self-consistency~\cite{sc} samples $K$ reasoning paths and votes, treating disagreement as \emph{noise to average away} under a sampling budget fixed in advance. \sys{} reads the same disagreement as \emph{information about the query}. When two samples agree, the query is easy for the current model and grounding, and a third sample is unlikely to change the outcome. Agreement thus becomes a stopping rule, not a voting mechanism. Adaptive-Consistency~\cite{adaptiveconsistency} and early-stopping self-consistency~\cite{escsc} both halt sampling once drawn answers converge, but they operate on a single query in isolation and must draw at least two samples before learning anything. \sys{} knows from Eq.~(\ref{eq:classifier}) that a point lookup cannot benefit from a second sample at all, and it skips the extra sample entirely for the majority of the workload. Agreement-based stopping and class-based budgeting are therefore complementary, the first acting within a class and the second across classes, and Section~\ref{subsec:exp4} measures them separately.

%!TEX root = ../main.tex
\section{Experiments}
\label{sec:experiments}

We organize the evaluation around eight questions. Exp-1 asks whether \sys{} outperforms state-of-the-art memory systems across memory abilities. Exp-2 asks whether the advantage survives a change of backbone. Exp-3 asks which query patterns the advantage comes from. Exp-4 asks what the accuracy costs in LLM calls. Exp-5 asks how much each architectural layer contributes. Exp-6 asks whether the second-order memory improves a user over time. Exp-7 asks how accurate the zero-cost planner is and how sensitive the end-to-end result is to its two free parameters. Exp-8 asks how read-path cost behaves as history length grows.

\subsection{Experimental Setup}
\label{subsec:setup}

%!TEX root = ../main.tex
\begin{table}[t]
    \centering
    \small
    \caption{The eight benchmarks, grouped by the memory ability they stress.
      History length spans two orders of magnitude, which lets us separate gains
      that come from structure from gains that come from context length.}
    \label{tbl:benchmarks}
    \setlength{\tabcolsep}{4pt}
    \begin{tabular}{@{}llr@{}}
        \toprule
        \textbf{Benchmark} & \textbf{Memory ability} & \textbf{History} \\
        \midrule
        \locomo~\cite{locomo}         & Multi-session dialogue    & $\sim$9K \\
        \lmev~\cite{longmemeval}      & Entity-centric recall     & $\sim$115K \\
        \membench~\cite{membench}     & Typed access primitives   & $\sim$100K \\
        \addlinespace[2pt]
        \lmeo~\cite{longmemeval}      & Long-horizon retention    & $\sim$40K \\
        \lmes~\cite{longmemeval}      & Long-horizon retention    & $\sim$115K \\
        \lmem~\cite{longmemeval}      & Long-horizon retention    & $\sim$1.5M \\
        \addlinespace[2pt]
        \evomem~\cite{evomemory}      & Evolving user state       & $\sim$50K \\
        \dynamicmem~\cite{dynamicmem} & Evolving user state       & $\sim$127K \\
        \bottomrule
    \end{tabular}
\end{table}


\stitle{Datasets.}
We use the eight long-term memory benchmarks in Table~\ref{tbl:benchmarks}, chosen to separate three memory abilities that a single benchmark would confound. \locomo~\cite{locomo} and \lmev~\cite{longmemeval} test recall over multi-session dialogue. \lmeo, \lmes, and \lmem~\cite{longmemeval} share one query style and scale the surrounding history by nearly two orders of magnitude, which lets history length be varied as a controlled factor. \evomem~\cite{evomemory} and \dynamicmem~\cite{dynamicmem} issue many questions per user over an evolving state, the only setting in which per-user adaptation becomes observable. \membench~\cite{membench} labels every one of its $300$ queries with the query pattern it requires, and it is therefore the one benchmark on which a per-class decomposition is possible.

\stitle{Baselines.}
We compare against seven systems in two groups. \etitle{(1) Prompting Baselines.} \textbf{Full-Context}~\cite{longcontext} places the entire history in the context window, and \textbf{RAG}~\cite{rag} retrieves top-$k$ passages by dense similarity. \etitle{(2) Memory Systems.} \textbf{MemGPT}~\cite{memgpt} pages facts between context tiers. \textbf{Mem0}~\cite{mem0} consolidates facts into a dictionary through an \textsc{Add}/\textsc{Update}/\textsc{Delete} loop, and \textbf{Mem0}$^{g}$ adds entity relations. \textbf{Zep}~\cite{zep} maintains a temporal knowledge graph. \textbf{LongMemEval-Fw}~\cite{longmemeval} applies a three-stage indexing, retrieval, and reading pipeline. All seven are run from their released implementations under the same backbone, which leaves the memory design as the sole source of any difference.

\stitle{Evaluation Metrics.}
Multiple-choice benchmarks are scored by exact match. The remaining benchmarks are scored by F1 against the gold answer, normalized at the phrase level for \lmev{} and at the token level elsewhere. Phrase-level normalization credits a prediction only when it reproduces the gold phrase, so absolute F1 on \lmev{} sits far below the token-level scores, and the two scales are not comparable; within \lmev{} the ranking across systems is unaffected. Higher is better throughout, and every reported number is an accuracy on the benchmark's own scale.

\stitle{Backbone LLMs and Settings.}
Unless stated otherwise all systems run on GPT-5.5. Exp-2 additionally evaluates three further closed-source backbones~\cite{gpt4,claude,gemini} and four open-source backbones, selected to span a wide capability range at two deployment tiers. \sys{} is implemented in Python against an OpenAI-compatible client. Every hyperparameter is fixed once and reused across all benchmarks and backbones, with no per-dataset tuning.

%!TEX root = ../main.tex
\begin{table*}[t]
    \centering
    \small
    \caption{Accuracy across eight long-term memory benchmarks. \sys{} matches
      or improves the strongest prior system on every benchmark, and the margin
      is widest on the long-horizon and state-inference groups where a single
      retrieval pass is least sufficient.}
    \label{tbl:main_results}
    \begin{tabular}{@{}l ccc ccc cc@{}}
        \toprule
        \multirow{2}{*}{\textbf{Method}}
            & \multicolumn{3}{c}{\textbf{Conversational memory}}
            & \multicolumn{3}{c}{\textbf{Long-horizon memory}}
            & \multicolumn{2}{c}{\textbf{Personalized state}} \\
        \cmidrule(lr){2-4} \cmidrule(lr){5-7} \cmidrule(l){8-9}
            & \textbf{\locomo} & \textbf{\lmev} & \textbf{\membench}
            & \textbf{\lmeo} & \textbf{\lmes} & \textbf{\lmem}
            & \textbf{\evomem} & \textbf{\dynamicmem} \\
        \midrule
        \tblgroup{9}{Prompting baselines} \\
        Full-Context~\cite{longcontext}   & 52.3 & 4.5 & 86.7 & 49.3 & 43.6 & 58.5 & 16.7 & 33.5 \\
        RAG~\cite{rag}                    & 53.1 & 4.9 & 87.8 & 50.8 & 45.1 & 59.1 & 17.3 & 34.1 \\
        \addlinespace[2pt]
        \tblgroup{9}{Memory systems} \\
        MemGPT~\cite{memgpt}              & 54.5 & 6.3 & 89.7 & 54.1 & 48.4 & 61.1 & 18.0 & 34.3 \\
        Mem0~\cite{mem0}                  & 56.4 & 7.2 & \second{91.2} & 58.3 & 54.2 & 65.2 & 20.4 & 34.5 \\
        Mem0$^{g}$~\cite{mem0}            & 57.1 & 7.5 & 90.8 & 60.1 & 56.7 & 67.1 & 21.2 & 34.7 \\
        Zep~\cite{zep}                    & \second{58.2} & 8.1 & 90.4 & \second{62.2} & \second{59.3} & \second{70.3} & \second{22.1} & \second{34.9} \\
        LongMemEval-Fw~\cite{longmemeval} & 57.3 & \second{9.4} & 90.1 & 61.7 & 58.8 & 69.4 & 21.6 & 34.6 \\
        \addlinespace[2pt]
        \ourgroup{9}{Ours} \\
        \textbf{\sys{}}
            & \best{58.8} & \best{10.3} & \best{92.8}
            & \best{67.2} & \best{64.9} & \best{75.2}
            & \best{24.6} & \best{34.9} \\
        \bottomrule
    \end{tabular}
\end{table*}


\subsection{Exp-1. Accuracy Across Three Memory Abilities}
\label{subsec:exp1}

We compare \sys{} against all seven baselines on the eight benchmarks of Table~\ref{tbl:benchmarks}, grouped by the memory ability each one stresses. Table~\ref{tbl:main_results} gives the comparison. \sys{} matches or improves the strongest prior system on every benchmark, and the margin varies systematically with the memory ability the benchmark stresses.

\stitle{Long-Horizon Memory.}
The margin is widest on the three LongMemEval variants, where it reaches $4.9$ to $5.6$ points, and it does not shrink as history grows by nearly two orders of magnitude. The reason lies in what each system hands the model. A longer history enlarges the candidate set that similarity retrieval returns, so prior systems pass the model more passages and leave it to locate, order, and align the relevant facts inside the context window, the regime in which long-context models are known to undercount and misorder~\cite{longcontext,ruler}. \sys{} reads the fields a query needs from the catalog indexes, so the context block stays small and already structured, and the model only performs the reasoning step. The size of that block is set by the query pattern, independent of $|\mathcal{H}|$, which is why the margin persists at the largest scale.

\stitle{Conversational Memory.}
The margin is narrowest on this group, at $0.6$ points on \locomo{} and $1.6$ on \membench{}. These benchmarks ask mostly for facts stated in one place, a query pattern that a single retrieval pass already serves well, so there is less headroom for an operator to recover. The result on \lmev{} is $10.3$ against $9.4$ for the strongest baseline, a $9.6\%$ relative improvement on a strict phrase-level F1 scale on which every system scores in the single digits. What distinguishes \sys{} on this group is that the accuracy and the cost advantage hold together: on \membench{} it reaches the highest accuracy while issuing fewer LLM calls than the uniform reference, as Exp-4 shows.

\stitle{Personalized State.}
The two state-inference benchmarks separate the systems least in absolute terms, with all seven baselines falling within $5.4$ points on \evomem{} and $1.4$ points on \dynamicmem. These benchmarks issue many questions per user, and a system that treats each query independently answers the tenth exactly as it answered the first, which is what prior designs do: they record what arrives on the write path and never what succeeded on the read path. An aggregate accuracy therefore averages over query position, the dimension on which the second-order memory acts, and hides any per-user improvement. Exp-6 conditions on query position and recovers the effect that the aggregate hides.

\finding{\sys{} matches or improves the strongest prior system on all eight benchmarks, and the margin is largest where a single retrieval pass is least sufficient, on long histories and on repeated questions from one user.}

\subsection{Exp-2. Backbone Generalization}
\label{subsec:exp2}

A system that supplies external structure should benefit every backbone, and the benefit should be largest where the backbone's own capacity is most limited. To test this, we hold \membench{} and the full \sys{} configuration fixed and vary only the backbone, over four closed-source and four open-source models spanning a wide capability range. Table~\ref{tbl:main_results_open} collects the comparison.

%!TEX root = ../main.tex
\begin{table}[t]
    \centering
    \small
    \caption{Accuracy on \membench{} across backbones. \textbf{Best prior} is
      the strongest of the seven baselines of Table~\ref{tbl:main_results}
      re-run on that same backbone. The gain is smallest where the backbone is
      strongest, so the structure \sys{} supplies substitutes for capacity the
      model does not have.}
    \label{tbl:main_results_open}
    \begin{tabular}{@{}lcc@{}}
        \toprule
        \textbf{Backbone} & \textbf{Best prior} & \textbf{\sys{}} \\
        \midrule
        \tblgroup{3}{Closed-source} \\
        GPT-5.5             & 91.2 & \best{92.8} \\
        Claude-3.5-Sonnet   & 87.4 & \best{91.4} \\
        Gemini-1.5-Pro      & 85.6 & \best{89.7} \\
        GPT-4o              & 82.8 & \best{88.7} \\
        \addlinespace[2pt]
        \tblgroup{3}{Open-source} \\
        DeepSeek-V4         & 82.7 & \best{88.5} \\
        Llama-3.1-70B       & 79.8 & \best{86.5} \\
        Qwen3-14B           & 77.4 & \best{85.1} \\
        Qwen3-8B            & 73.4 & \best{81.3} \\
        \bottomrule
    \end{tabular}
\end{table}


\sys{} improves every backbone, which rules out an explanation resting on one model's inductive bias. The size of the improvement is what varies. It is smallest on the strongest backbone and largest on the weakest, and this ordering holds within each group separately. This pattern suggests that external structure compensates for capability the backbone lacks. A frontier model already holds more history in context and counts over it more reliably, so the incremental gain is smaller, but the gain remains positive even on the strongest backbone. The practical consequence is that \sys{} widens the set of backbones on which long-term memory is usable, which matters most where frontier inference per query is not affordable.

\finding{The advantage of \sys{} is not backbone-specific, and it grows as backbone capability falls, with external memory structure substituting for internal model capacity.}

\subsection{Exp-3. Per-Pattern Accuracy Decomposition}
\label{subsec:exp3}

%!TEX root = ../main.tex
\begin{table}[t]
    \centering
    \small
    \caption{Accuracy by plan class on \membench. Every system handles point
      lookups well, and the systems separate on aggregation and comparison,
      which is where an explicit plan replaces free-form reasoning.}
    \label{tbl:perclass}
    \begin{tabular}{@{}lccc@{}}
        \toprule
        \textbf{Method} & \opPOINT & \opAGG & \opCMP \\
        \midrule
        \tblgroup{4}{Prompting baselines} \\
        Full-Context~\cite{longcontext}    & 90.9 & 64.4 & 84.2 \\
        RAG~\cite{rag}                     & 91.8 & 67.1 & 85.8 \\
        \addlinespace[2pt]
        \tblgroup{4}{Memory systems} \\
        MemGPT~\cite{memgpt}               & 92.9 & 68.5 & 87.0 \\
        Mem0~\cite{mem0}                   & 93.9 & 70.2 & \second{88.1} \\
        Mem0$^{g}$~\cite{mem0}             & 94.1 & \second{71.9} & 85.4 \\
        Zep~\cite{zep}                     & \second{94.3} & 71.2 & 85.4 \\
        LongMemEval-Fw~\cite{longmemeval}  & 93.8 & 70.8 & 84.7 \\
        \addlinespace[2pt]
        \ourgroup{4}{Ours} \\
        \textbf{\sys{}} & \best{94.6} & \best{78.0} & \best{91.2} \\
        \bottomrule
    \end{tabular}
\end{table}


Exp-1 established that \sys{} leads overall. To locate the source of that lead, we decompose \membench{} accuracy by the pattern each query requires, which \membench{} labels directly. Table~\ref{tbl:perclass} breaks the accuracy out by query pattern.

The decomposition reveals where the advantage originates. \sys{} leads on all three query patterns, and the gains concentrate on aggregation and comparison, where a specialized operator replaces unstructured reasoning. The separation appears on these two classes because they require the model to compose evidence across multiple records.

\stitle{Aggregation.}
Counting fails for reasons that a better retriever cannot fix. When candidate items are scattered through a long context, models miscount even after retrieving every relevant item~\cite{ruler}. Prior systems retrieve passages and ask for a count, which leaves the count implicit inside unstructured generation. \sys{} makes it explicit, enumerating candidates, deciding each one against the predicate, then tallying. Each step is bounded, and the total rests on enumerable evidence with an exact count.

\stitle{Comparison.}
Comparison fails through misalignment. The model must locate the compared attribute in two separate passages and align them before it can order them, and that alignment is what breaks on long histories. \sys{} performs the alignment outside the model, reading the attribute from the typed index for each candidate and placing both values adjacent in the prompt. The remaining task is an ordering over two values.

Both failure modes live in how the model \emph{processes} retrieved evidence, which is why improving retrieval alone does not close the gap and why the computation has to follow from the query pattern.

%!TEX root = ../main.tex
\begin{table}[t]
    \centering
    \small
    \caption{One aggregation query traced through three systems. The two
      baselines fail differently, and each failure is removed by a different
      layer of \sys{}.}
    \label{tbl:casestudy}
    \begin{tabular}{@{}llcc@{}}
        \toprule
        \textbf{System} & \textbf{Retrieval and reasoning} & \textbf{Calls} & \textbf{Answer} \\
        \midrule
        Full-Context   & count left implicit & 1 & \textcolor{dropred}{5} \\
        Zep~\cite{zep} & predicate unresolved & 1 & \textcolor{dropred}{4} \\
        \midrule
        \ourgroup{4}{\sys{}, layer by layer} \\
        \multicolumn{4}{@{}l}{\quad \emph{Planner} routes on the quantifier marker} \\
        \multicolumn{4}{@{}l}{\quad \emph{Second-order memory} fronts this user's resolved entities} \\
        \multicolumn{4}{@{}l}{\quad \emph{Operator} decides each candidate, then tallies} \\
        \multicolumn{4}{@{}l}{\quad \emph{Executor} withholds the third call on agreement} \\
        \textbf{\sys{}} & explicit per-item decision & 2 & \textcolor{cadmiumgreen}{\textbf{3}} \\
        \bottomrule
    \end{tabular}
\end{table}


To make the two failure modes concrete, Table~\ref{tbl:casestudy} traces one aggregation query through \sys{} and two baselines. Both baselines had the relevant material available, and in neither case did retrieval fail. The prompting baseline miscounts evidence spread across sessions, and the graph memory retrieves topically related nodes without resolving the predicate the query asks about. \sys{} removes the two failures with different layers. The planner selects the \agg{} pattern through Eq.~(\ref{eq:classifier}), the second-order memory promotes the entities this user's earlier queries resolved, the operator records an explicit decision per candidate before tallying, and the executor withholds its third call once two samples agree.

\finding{The accuracy advantage concentrates on aggregation and comparison, where a specialized operator replaces unstructured reasoning over scattered evidence.}

%!TEX root = ../main.tex
\begin{table}[t]
    \centering
    \small
    \caption{Calls per query on \membench. The two \sys{} rows share an accuracy
      column, which isolates the short-circuit as a pure cost mechanism and
      attributes the accuracy to the plan and not to the budget.}
    \label{tbl:cost}
    \setlength{\tabcolsep}{3pt}
    \begin{tabular}{@{}lcccrc@{}}
        \toprule
        \multirow{2}{*}{\textbf{Configuration}}
            & \multicolumn{3}{c}{\textbf{Calls per query}}
            & \multirow{2}{*}{\textbf{Total}} & \multirow{2}{*}{\textbf{Acc.}} \\
        \cmidrule(lr){2-4}
            & \opPOINT & \opAGG & \opCMP & & \\
        \midrule
        Uniform, one call            & 1.00 & 1.00 & 1.00 & 300 & 86.7 \\
        Uniform, three calls         & 3.00 & 3.00 & 3.00 & 900 & 89.1 \\
        \addlinespace[2pt]
        \sys{}, budget only          & 1.00 & 3.00 & 3.00 & 532 & 92.8 \\
        \textbf{\sys{}}              & 1.00 & 2.30 & 2.18 & \best{444} & \best{92.8} \\
        \bottomrule
    \end{tabular}
\end{table}

%!TEX root = ../../main.tex
% Cost-accuracy Pareto plot on MemBench.
\begin{figure}[t]
    \centering
    \resizebox{0.98\columnwidth}{!}{%
    \begin{tikzpicture}[
        every node/.style={font=\footnotesize},
        axis/.style={-Latex, thick, black},
        tick/.style={black, thin},
        grid/.style={gray!25, dashed, very thin},
        pt/.style={circle, minimum size=6pt, inner sep=0pt},
        lead/.style={gray!55, thin},
    ]
        \def\W{6.4}
        \def\H{3.6}
        %-------- axes with arrowheads only --------%
        \draw[axis] (0,0) -- (\W+0.15,0);
        \draw[axis] (0,0) -- (0,\H+0.15);
        %-------- x ticks --------%
        \foreach \x/\lbl in {0.0/1.0, 1.4/1.5, 2.8/2.0, 4.2/2.5, 5.6/3.0}{
            \draw[tick] (\x,0) -- (\x,-0.09);
            \node[below] at (\x,-0.09) {\lbl};
        }
        %-------- y ticks --------%
        \foreach \y/\lbl in {0.0/86, 0.9/88, 1.8/90, 2.7/92, 3.6/94}{
            \draw[tick] (0,\y) -- (-0.09,\y);
            \node[left] at (-0.09,\y) {\lbl};
        }
        %-------- axis labels --------%
        \node[below, font=\footnotesize] at (\W/2,-0.7) {\emph{Average LLM calls per query}};
        \node[rotate=90, font=\footnotesize] at (-0.9,\H/2) {\emph{Accuracy on \membench{} (\%)}};
        %-------- soft grid --------%
        \foreach \x in {1.4, 2.8, 4.2, 5.6} \draw[grid] (\x,0) -- (\x,\H);
        \foreach \y in {0.9, 1.8, 2.7, 3.6} \draw[grid] (0,\y) -- (\W,\y);
        %-------- data points --------%
        % x = (calls-1.0)*2.8;  y = (acc-86)*0.45
        \node[pt, fill=gray!70]    at ({(1.00-1.0)*2.8}, {(86.7-86)*0.45}) (fc)   {};
        \node[pt, fill=gray!70]    at ({(1.00-1.0)*2.8}, {(87.8-86)*0.45}) (rag)  {};
        \node[pt, fill=softblue!70] at ({(1.4-1.0)*2.8}, {(89.7-86)*0.45}) (mgpt) {};
        \node[pt, fill=softblue!70] at ({(1.5-1.0)*2.8}, {(91.2-86)*0.45}) (m0)   {};
        \node[pt, fill=softblue!70] at ({(1.8-1.0)*2.8}, {(90.8-86)*0.45}) (m0g)  {};
        \node[pt, fill=softblue!70] at ({(1.7-1.0)*2.8}, {(90.4-86)*0.45}) (zep)  {};
        \node[pt, fill=softblue!70] at ({(1.6-1.0)*2.8}, {(90.1-86)*0.45}) (lfw)  {};
        \node[pt, fill=softorange] at ({(3.00-1.0)*2.8}, {(89.1-86)*0.45}) (uni)  {};
        \node[pt, fill=dropred, minimum size=8pt] at ({(1.48-1.0)*2.8},{(92.8-86)*0.45}) (mc) {};
        %-------- pareto arrow --------%
        \draw[dropred, dashed, thick, opacity=0.5] (fc.center) -- (mc.center);
        %-------- labels with leaders, dispersed --------%
        \node[anchor=west] (mcL) at (2.3, 3.35) {\textbf{\sys{} (ours)}};
        \draw[lead] (mc) -- (mcL.west);
        \node[anchor=west] (m0L)  at (2.5, 2.60) {Mem0 (91.2)};
        \draw[lead] (m0) -- (m0L.west);
        \node[anchor=west] (m0gL) at (2.9, 2.15) {Mem0$^{g}$ (90.8)};
        \draw[lead] (m0g) -- (m0gL.west);
        \node[anchor=west] (zepL) at (2.6, 1.80) {Zep (90.4)};
        \draw[lead] (zep) -- (zepL.west);
        \node[anchor=west] (lfwL) at (2.3, 1.45) {LongMemEval-Fw (90.1)};
        \draw[lead] (lfw) -- (lfwL.west);
        \node[anchor=west] (mgptL) at (2.0, 1.10) {MemGPT (89.7)};
        \draw[lead] (mgpt) -- (mgptL.west);
        \node[anchor=east] (uniL) at (5.5, 0.60) {Uniform-$K{=}3$ (89.1)};
        \draw[lead] (uni) -- (uniL.east);
        \node[anchor=west] (ragL) at (0.25, 0.90) {RAG (87.8)};
        \draw[lead] (rag) -- (ragL.west);
        \node[anchor=west] (fcL) at (0.25, 0.30) {Full-Context (86.7)};
        \draw[lead] (fc) -- (fcL.west);
    \end{tikzpicture}
    }
    \caption{Cost--accuracy Pareto on \membench{} (GPT-5.5, $n{=}300$). Each point is a system's average LLM calls per query (x-axis) versus its accuracy (y-axis). \sys{} occupies the upper-left corner: strictly higher accuracy than every prior memory system and strictly lower cost than the Uniform-$K{=}3$ reference. No prior system Pareto-dominates \sys{}.}
    \Description{Scatter plot of average LLM calls per query against accuracy on
      MemBench. MemCatalog sits in the upper-left region, with higher accuracy
      than every prior memory system and lower cost than the uniform K equals
      three reference.}
    \label{fig:pareto}
\end{figure}


\subsection{Exp-4. Cost Analysis}
\label{subsec:exp4}

Accuracy bought with extra LLM calls is an artifact of the budget, and the architecture should be evaluated separately. To separate the two, we measure the LLM calls each system issues per query on \membench{} and place cost and accuracy on the same axes. Cost is counted on the read path for every system, so a system that spends a retrieval or rewriting call before generating is charged for it. Write-path extraction is a one-time per-session cost that every structured system amortizes over the queries touching that session, including the on-demand extraction of \sys{}, and it is excluded uniformly for all systems. We additionally include two reference configurations that fix the sampling budget uniformly at one call and at three calls per query, which isolates the effect of \emph{allocating} a budget from the effect of its \emph{size}.

Figure~\ref{fig:pareto} gives the result. \sys{} occupies the upper-left region, where no evaluated system attains both lower cost and comparable accuracy. The two groups miss it for different reasons. Prior memory systems are inexpensive and lose accuracy on the queries a single pass cannot answer. The uniform three-call reference recovers part of that accuracy and remains dominated, because it spends its extra samples on the point lookups that never required them.

Table~\ref{tbl:cost} separates the two mechanisms behind this position. Pattern-aware budgeting acts across queries, assigning the single-call budget to the point lookups that form the majority of the workload. Agreement-based early stopping, the rule of Eq.~(\ref{eq:short}), acts within the remaining queries and stops early on those whose two samples already agree. The first mechanism draws from the easy pattern and the second from the easy instances of the hard patterns, and the slices are disjoint, which lets the savings compose. The accuracy column settles what the savings cost: it is unchanged between the two \sys{} configurations, so early stopping removes $88$ of $532$ calls without giving up a single answer, and the accuracy of \sys{} is attributable to the framework.

\finding{\sys{} is Pareto-optimal against every evaluated system, and early stopping lowers its cost at unchanged accuracy, which attributes the accuracy to the framework, with early stopping purely reducing cost.}

%!TEX root = ../main.tex
\begin{table}[t]
    \centering
    \small
    \caption{Removing one component at a time. Each layer contributes on a
      distinct axis, accuracy for the framework, accuracy under repeated
      queries for the second-order memory, and cost for the executor,
      confirming the three are complementary.}
    \label{tbl:ablation}
    \begin{tabular}{@{}lccccc@{}}
        \toprule
        \multirow{2}{*}{\textbf{Configuration}}
            & \multicolumn{2}{c}{\textbf{\membench}}
            & \multicolumn{2}{c}{\textbf{\lmev}} \\
        \cmidrule(lr){2-3} \cmidrule(l){4-5}
            & \textbf{Acc.} & \textbf{Calls} & \textbf{Acc.} & \textbf{Calls} \\
        \midrule
        \textbf{\sys{}}               & \best{92.8} & \best{444} & \best{10.3} & \best{602} \\
        \addlinespace[2pt]
        \tblgroup{5}{Framework} \\
        w/o planner          & 89.1 & 900 & 8.8 & 903 \\
        w/o physical operators        & 87.9 & 900 & 8.3 & 903 \\
        w/o catalog                   & 86.7 & 900 & 7.6 & 903 \\
        \addlinespace[2pt]
        \tblgroup{5}{Second-order memory} \\
        w/o workload log              & 92.1 & 444 & 8.6 & 602 \\
        w/o index warming             & 92.5 & 444 & 9.3 & 602 \\
        w/o session eviction          & 92.8 & 444 & 10.1 & 651 \\
        \addlinespace[2pt]
        \tblgroup{5}{Cost-aware executor} \\
        w/o plan-class budget         & 92.5 & 756 & 10.3 & 903 \\
        w/o stability short-circuit   & 92.8 & 532 & 10.3 & 723 \\
        \bottomrule
    \end{tabular}
\end{table}


\subsection{Exp-5. Ablation Study}
\label{subsec:exp5}

The three layers of \sys{} were motivated by three different obstacles. This experiment confirms that each layer contributes independently by removing one component at a time and measuring both accuracy and cost on \membench{} and \lmev{}, a multiple-choice benchmark and an open-answer one. Table~\ref{tbl:ablation} collects the ablations.

Each layer contributes on a distinct axis, confirming they are complementary. Removing any framework component costs accuracy and simultaneously removes the pattern that the budget depends on, so cost rises to the uniform setting. Within the framework the ordering is informative. Removing the catalog is worse than removing the planner, since a planner that routes correctly still has nothing typed to read. The second-order memory contributes most on \lmev{}, where a user's repeated questions benefit from read-path feedback, and its three components are not equally active on both benchmarks: session eviction leaves \membench{} untouched, whose histories stay within the cache, while on \lmev{} it saves $49$ calls. In the executor the two mechanisms differ in kind. Pattern-aware budgeting lowers cost on both benchmarks and also contributes $0.3$ points on \membench, since a point lookup that is sampled three times can be outvoted on the one answer the index already supplies. Agreement-based early stopping changes no accuracy figure on either benchmark and is a cost mechanism alone.

\finding{Each layer contributes on a distinct axis, accuracy for the framework, accuracy under repeated queries for the second-order memory, and cost for the executor, confirming they are complementary.}

%!TEX root = ../main.tex
\begin{table}[t]
    \centering
    \small
    \caption{Accuracy against the position of a query in a user's own sequence,
      evaluated on users with $\geq 10$ queries so that all four buckets are
      populated. Only \sys{} improves as a user keeps asking. On \dynamicmem{}
      it starts below the static rankers in the first bucket, where its
      workload log is still empty, and overtakes them once the log carries
      signal. Absolute levels differ from Table~\ref{tbl:main_results}, which
      averages over all users.}
    \label{tbl:second_order_learning}
    \begin{tabular}{@{}lcccc@{}}
        \toprule
        \multirow{2}{*}{\textbf{Method}}
            & \multicolumn{4}{c}{\textbf{Query position for one user}} \\
        \cmidrule(l){2-5}
            & 1--3 & 4--6 & 7--9 & 10+ \\
        \midrule
        \tblgroup{5}{\evomem} \\
        Full-Context                   & 16.7 & 17.0 & 16.8 & 16.5 \\
        Mem0~\cite{mem0}               & 20.1 & 20.5 & 20.4 & 20.2 \\
        Zep~\cite{zep}                 & \second{21.8} & \second{22.2} & \second{22.1} & \second{22.0} \\
        \textbf{\sys{}}                & \best{22.2} & \best{24.4} & \best{26.1} & \best{27.8} \\
        \addlinespace[2pt]
        \tblgroup{5}{\dynamicmem} \\
        Full-Context                   & 28.7 & 28.8 & 28.6 & 28.4 \\
        Mem0~\cite{mem0}               & \second{29.4} & 29.8 & 29.6 & 29.3 \\
        Zep~\cite{zep}                 & \best{29.9} & \second{30.0} & \second{29.9} & \second{29.8} \\
        \textbf{\sys{}}                & 29.1 & \best{30.3} & \best{31.7} & \best{32.8} \\
        \bottomrule
    \end{tabular}
\end{table}


\subsection{Exp-6. Accuracy by Query Position}
\label{subsec:exp6}

The second-order memory predicts a specific and falsifiable behavior. Accuracy on one user should rise as that user asks more questions, and no system without a read-path feedback loop should show the same rise. To test this we restrict to users with at least ten queries, so that all four position buckets are populated, condition on the position of a query within its own user's sequence, and report accuracy per position bucket, holding backbone, prompt, and benchmark fixed so that position is the only variable. Table~\ref{tbl:second_order_learning} bins the accuracy by position. Because the evaluation population is a subset of the full test set, the absolute accuracy levels in each column differ from Table~\ref{tbl:main_results}; the within-table ranking and the monotone trend are the subject of this experiment.

The predicted separation appears on both benchmarks. \sys{} improves monotonically across position buckets, gaining $5.6$ points on \evomem{} and $3.7$ on \dynamicmem{} from the first bucket to the last, while no baseline moves by more than $0.5$ points in either direction and none shows a monotone trend, since none of them records which stored entries produced accepted answers and none has a signal that could vary with position. The starting point is equally informative. In the first bucket the workload log is still empty, so \sys{} has no second-order signal yet and begins level with the static rankers, below Zep by $0.8$ points on \dynamicmem. It overtakes them from the second bucket onward and leads by $3.0$ points by the last. A mechanism that learns from answered queries should hold no advantage before any query has been answered, and this is the crossover that behavior produces.

The improvement cannot be attributed to the model or the prompt, both of which are identical across buckets. It follows the priority score of Eq.~(\ref{eq:pi}), which accumulates evidence that an entity participates in answers this user accepts. By the later buckets the context block handed to the operator is composed largely of entities with a demonstrated history of being queried, where a static ranker still offers entities selected by similarity to the current query alone. This is the read-path signal of Section~\ref{sec:secondorder}, and Exp-5 confirms that removing the workload log removes the effect.

\finding{Only \sys{} improves as a user continues asking, which isolates read-path feedback as the mechanism and shows that prior systems are structurally unable to exhibit this behavior.}

\subsection{Exp-7. Planner Accuracy and Parameter Stability}
\label{subsec:exp7routing}

%!TEX root = ../main.tex
\begin{table}[t]
    \centering
    \small
    \caption{Query-classification quality and parameter stability on \membench. The
      upper block compares the rule-based planner against an LLM planner
      that spends one call per query and against a fixed planner. The lower
      block varies the two parameters \sys{} fixes system-wide around their
      defaults $\theta{=}0.6$ and $\alpha{=}0.5$, holding all else equal.}
    \label{tbl:routing}
    \setlength{\tabcolsep}{4pt}
    \begin{tabular}{@{}lccc@{}}
        \toprule
        \textbf{Planner} & \textbf{Class acc.} & \textbf{Planner calls} & \textbf{End-to-end acc.} \\
        \midrule
        Rule-based (ours) & 91.4 & 0 & 92.8 \\
        LLM planner          & 95.1 & 1 & 93.2 \\
        Always \point        & 61.3 & 0 & 87.5 \\
        \midrule
        \tblgroup{4}{Parameter perturbation} \\
        $\theta \in [0.4, 0.8]$  & \multicolumn{3}{c}{end-to-end acc.\ 92.3--92.9} \\
        $\alpha \in [0.3, 0.7]$  & \multicolumn{3}{c}{end-to-end acc.\ 92.2--92.8} \\
        \bottomrule
    \end{tabular}
    \Description{Planner comparison and parameter sensitivity on MemBench.}
\end{table}


The planner governs which operator handles each query, so its quality directly affects the end-to-end result. This experiment validates its design along two dimensions, classification quality and parameter stability. We measure classification accuracy against the class labels \membench{} supplies, comparing the rule-based planner against an LLM planner that spends one call per query and against a fixed planner that assigns every query to \point. We then perturb the early-stopping threshold and the leading weight of the priority score over wide intervals and re-measure end-to-end accuracy. Table~\ref{tbl:routing} covers both.

Three observations follow from the table. First, the rule-based planner reaches $91.4\%$ classification accuracy at zero LLM cost. Second, classification quality has already saturated with respect to what the system does with it: an LLM planner classifies better, at $95.1\%$, yet buys only $+0.4$ end-to-end accuracy for one extra call on every query, which is $300$ additional calls against a total of $444$, and a gap of $0.4$ is no larger than the spread the parameter perturbations below produce. Paying two-thirds more calls for a difference of that size is not a trade this workload rewards. Third, the fixed planner locates where the value of query classification actually sits: routing every query to \point{} costs $5.3$ points end-to-end, an order of magnitude more than separates any two planner variants, so what matters is that the query pattern is chosen at all, and the choice granularity has diminishing returns. The lower block perturbs the two parameters \sys{} fixes system-wide, the early-stopping threshold $\theta$ and the leading weight $\alpha$ of the priority score, across intervals spanning half of each parameter's usable range. End-to-end accuracy stays within $0.6$ points over both, which leaves the reported result independent of any parameter search.

\finding{The rule-based planner reaches $91.4\%$ classification accuracy at zero cost, an LLM planner cannot convert better classification into materially better answers, and end-to-end accuracy is stable across wide parameter intervals.}

\subsection{Exp-8. Scalability Analysis}
\label{subsec:exp8}

%!TEX root = ../main.tex
\begin{table}[t]
    \centering
    \small
    \caption{Behavior as history grows on the three LongMemEval variants.
      Read-path cost is governed by the plan class and not by history length,
      so it stays flat while the prompting baseline grows with the context
      it must carry.}
    \label{tbl:scaling}
    \setlength{\tabcolsep}{4pt}
    \begin{tabular}{@{}lccc@{}}
        \toprule
        \textbf{History} & \textbf{Calls / query} & \textbf{Latency (s)} & \textbf{Accuracy} \\
        \midrule
        \tblgroup{4}{Full-Context} \\
        $\sim$40K  (\lmeo)  & 1.0 & 2.3  & 49.3 \\
        $\sim$115K (\lmes)  & 1.0 & 5.8  & 43.6 \\
        $\sim$1.5M (\lmem)  & 1.0 & 47.6 & 58.5 \\
        \addlinespace[2pt]
        \tblgroup{4}{\sys{}} \\
        $\sim$40K  (\lmeo)  & 1.9 & 3.2 & 67.2 \\
        $\sim$115K (\lmes)  & 1.9 & 3.4 & 64.9 \\
        $\sim$1.5M (\lmem)  & 1.9 & 3.5 & 75.2 \\
        \bottomrule
    \end{tabular}
    \Description{Cost and accuracy as history length grows from 40K to 1.5M tokens.}
\end{table}


A memory system is only interesting if its cost survives the growth that motivates it. The three LongMemEval variants share one query style and scale the surrounding history by nearly two orders of magnitude, from $40$K to $1.5$M tokens, which isolates history length as the varied factor. We measure calls per query, end-to-end latency, and accuracy at each scale against the prompting baseline, whose prompt necessarily carries the history it searches. Table~\ref{tbl:scaling} tracks all three.

Cost separates the two systems along the length axis. Calls per query are flat for both, at $1.9$ for \sys{} and $1.0$ for Full-Context, so latency is where the difference appears. Latency grows $20.7\times$ for Full-Context, from $2.3$s to $47.6$s, and $1.09\times$ for \sys{}, from $3.2$s to $3.5$s. At the shortest history \sys{} is the slower of the two, by $0.9$s, because catalog access is a fixed overhead that a single prompting call does not pay. That overhead does not grow, so it is already repaid at $115$K, where \sys{} answers in $3.4$s against $5.8$s, and at $1.5$M it is $13.6\times$ faster. What the table shows is a different asymptote, with \sys{} converging to a flat cost while the baseline grows linearly. Accuracy moves with the difficulty of each variant for both systems, and the ranking does not: \sys{} leads at all three scales, by $17.9$, $21.3$, and $16.7$ points, with no erosion as history grows.

The asymmetry is structural and follows from where each system spends. A prompting baseline must carry the history it searches, so its prompt grows with $|\mathcal{H}|$ and its latency grows with it. The read-path cost of \sys{} is set by the query pattern through Eq.~(\ref{eq:cost}), which depends on the shape of the query alone, independent of the size of the store, and adding sessions changes what the catalog contains without changing what a query costs to answer.

\finding{Read-path cost is governed by the query pattern, independent of history length, which turns the growth of the store from a linear cost into a constant one and leaves the accuracy advantage of \sys{} intact at every scale.}

%!TEX root = ../main.tex
\section{Related Work}
\label{sec:related}

\subsection{Agent Memory Systems and Benchmarks}

The dominant thread in agent memory research treats memory as a \emph{storage} abstraction. MemGPT~\cite{memgpt} and its framework successor Letta~\cite{letta} draw inspiration from operating-system memory hierarchies. They treat the LLM context window as a fast tier and external stores as slow tiers, with paging between them managed by the agent via function calls. Their contribution is the movement policy across tiers, and the query-time algebra over the stored data remains a single retrieval pass. Mem0~\cite{mem0} designs a write-time pipeline that extracts, evaluates, and consolidates salient facts through an \textsc{Add}/\textsc{Update}/\textsc{Delete}/\textsc{Noop} loop, and its graph variant Mem0$^{g}$ adds entity and relation edges. On the read side Mem0 falls back to dense retrieval top-$k$, without a query planner or any optimization specific to the query pattern. Zep~\cite{zep}, powered by the Graphiti engine, models memory as a temporal knowledge graph over episodes, entities, and communities. It is tuned for retrieval accuracy and latency, but every query still goes through the same retrieval strategy. Further work such as A-MEM~\cite{agenticmem}, AriGraph~\cite{arigraph}, HippoRAG~\cite{hipporag}, MemoRAG~\cite{memorag}, and hierarchical retrieval methods such as RAPTOR~\cite{raptor} and GraphRAG~\cite{graphrag} add autonomous memory evolution, biologically inspired retrieval, and global memory augmentation. Skill-library agents such as Voyager~\cite{voyager} treat learned skills as reusable memory. All of these systems enrich the storage layer, as chronicled by a recent survey~\cite{agentmemsurvey}. None provides a query planner over typed operators, a cost model that budgets LLM calls by query pattern, or a workload feedback loop. \sys{} fills these three gaps in Sections~\ref{sec:framework}, \ref{sec:executor}, and \ref{sec:secondorder}.

Long-term memory evaluation has matured in parallel. LoCoMo~\cite{locomo} introduces multi-session dialogue at scale, and LongMemEval~\cite{longmemeval} covers five memory abilities up to 1.5M tokens while contributing the three-stage pipeline our LongMemEval-Fw baseline follows. Further benchmarks include MemoryBank~\cite{memorybank}, PerLTQA~\cite{perltqa}, DialSim~\cite{dialsim}, and LongBench~\cite{longbench}.

\subsection{DBMS Foundations for LLM Systems}

Two complementary DBMS lines meet in this paper. First, a body of work applies DBMS abstractions to LLM analytics over \emph{typed tabular data}. LOTUS~\cite{lotus} defines the first declarative model for AI-based tabular transformations such as semantic filter, join, top-$k$, and group-by, with gold algorithms and cost-optimized execution plans that achieve up to $1{,}000\times$ speedup. DocETL~\cite{docetl}, UQE~\cite{uqe}, and SUQL~\cite{suql} address related settings for unstructured documents and hybrid structured/unstructured search. These works share the core intuition of \sys{}, namely to lift LLM operations to a queryable algebra, but their target data is structured (DataFrames, tables). Agent memory is unstructured, its schema is latent in natural language, and its queries are user questions expressed as free text. To our knowledge \sys{} is the first system to transfer the full query-processing stack (planner, operators, cost model, and statistics) to agent memory.

Second, three founding lines of RDBMS research supply the pillars \sys{} borrows, namely cost-driven access-path selection~\cite{selinger1979}, extensible operator libraries~\cite{graefe1993}, and workload-driven self-tuning~\cite{chaudhuri1997autotuning,chaudhuri2007selftuning} that learned optimizers~\cite{neoopt,baoopt,balsa} later extended. The planner, the operator library, and the second-order memory of \sys{} play these three roles in turn. The adaptation is non-trivial, since the operators reason over natural language, the records are unstructured, and the cost unit becomes an LLM call, but the separation of concerns transfers intact.

\subsection{LLM Reasoning Primitives}

Chain-of-Thought~\cite{cot}, Self-Consistency~\cite{sc}, ReAct~\cite{react}, Reflexion~\cite{reflexion}, and Tree-of-Thoughts~\cite{tot} form the reasoning substrate on which the operators of \sys{} build. A line of work reduces the cost of self-consistency by adapting how many samples a query receives, halting once drawn answers agree~\cite{adaptiveconsistency,escsc}. These methods decide within a single query and must draw at least two samples before any decision is possible. \sys{} composes with them from the opposite direction, using the query pattern to decide how much budget a query is eligible for before sampling begins, which removes the extra sample entirely for the class that cannot benefit from it (Section~\ref{sec:executor}).


\section{Conclusion}
\label{sec:conclusion}

We introduced \sys{}, a DBMS-style memory management system for LLM agents. \sys{} organizes agent memory as managed data with a catalog, a zero-cost planner, and three specialized operators, and augments this framework with two optimization layers: a \emph{Second-Order Memory} that lets storage adapt to per-user workload, and a \emph{Cost-Aware Executor} that allocates its call budget by query pattern and by sample agreement. On eight long-term memory benchmarks and eight LLM backbones, \sys{} matches or improves the strongest prior memory system on every benchmark, by up to $5.6$ points on \lmes. On \membench{} it further reduces LLM calls by $50.7\%$ relative to a Uniform-$K{=}3$ reference while attaining higher accuracy, and its second-order memory is the only mechanism among the evaluated systems whose accuracy rises monotonically across successive queries from the same user. The classical DBMS separation of storage, planning, and execution therefore transfers cleanly to agent memory, and each of the three layers contributes measurably. We view \sys{} as a first step toward treating LLM agent memory as first-class managed data, a setting in which decades of data management results become applicable.

Future work will explore multi-user catalog sharing under privacy constraints, tighter integration of \sys{} with retrieval-native architectures that pre-index by query pattern, and extensions of the operator library to multi-hop reasoning that currently falls back to \agg{} enumeration.


\begin{acks}
Anonymized for submission.
\end{acks}

\bibliographystyle{ACM-Reference-Format}
\bibliography{memcatalog}

\end{document}
